% page 297 du fac-simile (image p-296 du scan)
% bloc 160.0 x 233.3 mm ; marge gauche 42.0 mm
\pgc{20.320mm}{0.000mm}{
\l{\cel{52}289}
\l{}
\l{}
\l{komisar. (\sou{ulu, pri}). - Konfidar (ad ulu) la tasko qua konsistas}
\l{\cel{3}ye agar o facar (to o co). - DEFIS.}
\l{}
\l{komisario. La persono qua esas delegita pri ula funcioni,}
\l{\cel{3}pri ula ofico - provizore o longa-dure. DEFIRS.}
\l{}
\l{komisiono.\cel{2}To quon privato demandas de ulu ke ica agez vice}
\l{\cel{3}lu (kompro, quero, e c.) - DEFIRS.}
\l{}
\l{komisuro. (\sou{anat.}) Singla ek la du anguli ube la labii inter-}
\l{\cel{3}unionas. - DEFIS.}
\l{}
\l{komitato.\cel{2}Grupo de personi, selektita ek korpo plu personoza,}
\l{\cel{3}por divenar asemblitaro specala di qua la tasko esas preci-}
\l{\cel{3}pue exekutigala. - DEFIRS.}
\l{}
\l{komizo. Employato di administreyo, di komerco-firmo, di banko-}
\l{\cel{3}firmo. - DFI.}
\l{}
\l{komocionar. (\cel{1}\sou{trans}.) Shanceligar subite. - DEFIS.}
\l{}
\l{komoda. Quan karakterizas la uzo-facileso segun-deziro.}
\l{\cel{3}- DEFIS.}
\l{}
\l{komodo.\cel{2}Moblo, di qua la suprajo ne esas plu alta kam la}
\l{\cel{3}adulto-kubiti, garnisita ye tirkesti en qui onu lokizas}
\l{\cel{3}linjo, vesti, e c. - DEFIS.}
\l{}
\l{komodoro. I. (\sou{Navaro}\cel{1}Britaniana ed Usana). La oficiro di qua}
\l{\cel{3}la grado esas nemediate sub olta di sub-amiralo. - II.}
\l{\cel{3}(\sou{Navaro Nederlandana}). Navestro qua komandas eskadro. -}
\l{\cel{3}DEFIS.}
\l{}
\l{komono. I. (\sou{lor feudismo}). Urbo liberigita de la yugo feudala,}
\l{\cel{3}e di qua la borgezi su administris ipse, per oficistaro}
\l{\cel{3}elektita. - II. En Francia : urbo, burgo, vilajo o grupo}
\l{\cel{3}de domi, ye qua konsistas la unajo teritoriala quan adminis-}
\l{\cel{3}tras \textquotedbl{}urbestro\textquotedbl{} kun \textquotedbl{}+konselio municipala\textquotedbl{}. - III. En Bri-}
\l{\cel{3}tania : urbo, burgo, qua nominas deputiti a la Parlamento. -}
\l{\cel{3}- DEFIRS.}
\l{}
\l{kompakta.\cel{2}(\sou{aludante korpo solida}). Di qua la mas-uni esas tre}
\l{\cel{3}inter-kompresanta. - DEFIRS.}
\l{}
\l{kompano.- I. La persono qua vivas kustumale intima-relate kun}
\l{\cel{3}ulu. - II. (\sou{olim)}.\cel{2}Laboristo qua, ne-patrona, iris laborar}
\l{\cel{3}che patrono kun altra laboristi. - EFR.}
\l{}
\l{kompanio. I. (\sou{komerco})\cel{2}La ssociito o la asociiti, aludita da}
\l{\cel{3}la abreviuro \textquotedbl{}e Ko.\textquotedbl{} - II. Parto di bataliono, quan komandas}
\l{\cel{3}kapitano. -DEFIRS.}
}
